\documentclass[conference]{IEEEtran}
\IEEEoverridecommandlockouts
% Template version as of 6/27/2024

\usepackage{cite}
\usepackage{amsmath,amssymb,amsfonts}
\usepackage{algorithmic}
\usepackage{graphicx}
\usepackage{textcomp}
\usepackage{xcolor}
\usepackage[utf8]{inputenc}
\usepackage[T1]{fontenc}
\usepackage[brazil]{babel}
\usepackage{booktabs}
\usepackage{multirow}
\usepackage{url}

\def\BibTeX{{\rm B\kern-.05em{\sc i\kern-.025em b}\kern-.08em
    T\kern-.1667em\lower.7ex\hbox{E}\kern-.125emX}}

\begin{document}

\title{Análise Estatística de Métricas de Código-Fonte do Framework Neo4j}

\author{\IEEEauthorblockN{1\textsuperscript{st} Igor Mariano Lopes Rodrigues}
\IEEEauthorblockA{\textit{Curso de Engenharia de Software} \\
\textit{IBMEC}\\
Rio de Janeiro, Brasil \\
202407095992@alunos.ibmec.edu.br}
\and
\IEEEauthorblockN{2\textsuperscript{nd} Filipe Oliveira de Saldanha da Gama}
\IEEEauthorblockA{\textit{Curso de Engenharia de Software} \\
\textit{IBMEC}\\
Rio de Janeiro, Brasil \\
202303161085@alunos.ibmec.edu.br}
}

\maketitle

\begin{abstract}
This paper presents a statistical analysis of source code metrics from the Neo4j
framework using the GitHub Bug Dataset v1.1, covering class-level (7,917 instances)
and file-level (4,261 instances) granularities. Descriptive statistics (central
tendency, dispersion, quartiles), normality tests (Shapiro-Wilk and Lilliefors),
Pearson and Spearman correlations, VIF multicollinearity analysis, cost-sensitive
logistic regression, and 10-fold cross-validation were applied. Strong correlations
among size and complexity metrics were found (e.g., LOC x LLOC: r = 0.99).
All distributions are non-normal, and a severe class imbalance (less than 0.15 percent positive)
fundamentally limits predictive modeling: with only 8 buggy classes and 6 buggy
files, the Events Per Variable ratio is far below the recommended minimum of 10,
rendering any classifier unreliable. The main finding is dataset inadequacy for
fault prediction, not the inadequacy of the metrics themselves.
\end{abstract}

\begin{IEEEkeywords}
software metrics, object-oriented design, statistical analysis, fault prediction, class imbalance, Neo4j
\end{IEEEkeywords}

% =====================================================================
\section{Introdução}\label{sec:intro}

A análise quantitativa de métricas de código-fonte constitui abordagem fundamental na Engenharia de Software para compreender qualidade, complexidade e propensão a defeitos~\cite{pressman2014}. As métricas orientadas a objetos propostas por Chidamber e Kemerer~\cite{chidamber1994} permitem avaliar características estruturais de projetos Java e são amplamente utilizadas em predição de falhas~\cite{gyimothy2005,basili1996}.

O framework Neo4j é um sistema de banco de dados em grafo de código aberto, amplamente adotado em aplicações que demandam modelagem de relacionamentos complexos. Por sua relevância no ecossistema Java e pelo volume considerável de código-fonte, constitui um caso de estudo relevante para análise de métricas de software.

Este trabalho aplica técnicas de estatística descritiva e inferencial sobre datasets de métricas do framework Neo4j, extraídos do GitHub Bug Dataset v1.1~\cite{toth2016}. A análise contempla: (i)~medidas de tendência central, dispersão e posição relativa; (ii)~gráficos e detecção de outliers; (iii)~testes de normalidade (Shapiro-Wilk e Lilliefors); (iv)~correlação (Pearson e Spearman) e multicolinearidade (VIF); e (v)~regressão logística \textit{cost-sensitive} com validação cruzada. O repositório está disponível em \url{https://github.com/IgorMariano25/Projeto-ES-Neo4j}.

% =====================================================================
\section{Metodologia}\label{sec:metodo}

\subsection{Fonte de Dados}

Os datasets foram obtidos do GitHub Bug Dataset v1.1~\cite{toth2016}. Para o Neo4j, foram selecionados os datasets mais recentes: \textbf{Class-level} com 7.917 instâncias e 113 atributos (complexidade, acoplamento, coesão, herança, tamanho, documentação, \textit{warnings}) e \textbf{File-level} com 4.261 instâncias e 12 atributos (complexidade ciclomática McCC, LLOC e histórico de modificações). Nenhum dataset apresentou valores ausentes.

\subsection{Seleção de Variáveis}

A seleção aplicou três critérios objetivos sobre os 113 atributos Class-level e os 12 atributos File-level. \textbf{(i) Natureza métrica}: foram retidas apenas variáveis com semântica de propriedade quantificável do código-fonte; foram descartados identificadores textuais (\texttt{ID}, \texttt{Name}, \texttt{LongName}, \texttt{Parent}, \texttt{Component}, \texttt{Path}) e metadados posicionais (\texttt{Line}, \texttt{Column}, \texttt{EndLine}, \texttt{EndColumn}). \textbf{(ii) Variabilidade não-nula}: variáveis constantes foram excluídas, pois invalidam correlações, testes de hipótese e o cálculo do VIF; no File-level esse critério eliminou \textbf{CLOC} (\textit{Comment Lines of Code}), cuja totalidade das 4.261 instâncias é igual a zero. \textbf{(iii) Não-redundância}: no Class-level, foram descartadas as variantes com prefixo \textbf{T} (totais agregados a classes aninhadas, com $r > 0{,}99$ contra suas contrapartes), as variantes locais (prefixo \texttt{NL}) e \textit{proxies} redundantes (CBOI, NII, NOI) já capturados por CBO/RFC; as 37 categorias granulares de regras PMD foram colapsadas nos agregados \texttt{WarningInfo} e \texttt{WarningMajor}.

Restam \textbf{19 variáveis numéricas no Class-level} (Tabela~\ref{tab:selecao}) e \textbf{7 no File-level}: McCC, LLOC, \textit{Number of previous modifications}, \textit{Number of previous fixes}, \textit{Number of committers}, \textit{Number of developer commits} e a variável-alvo \textit{Number of bugs}.

\begin{table}[htbp]
\caption{Variáveis selecionadas para análise (Class-level)}
\label{tab:selecao}
\centering
\scriptsize
\setlength{\tabcolsep}{3pt}
\renewcommand{\arraystretch}{1.15}
\begin{tabular}{@{}p{0.30\columnwidth}p{0.20\columnwidth}p{0.42\columnwidth}@{}}
\toprule
\textbf{Variável} & \textbf{Categoria} & \textbf{Descrição} \\
\midrule
WMC, NL, NLE              & Complexidade  & Métodos ponderados; aninhamento \\
CBO, RFC                  & Acoplamento   & Acoplamento entre objetos; RFC \\
LCOM5                     & Coesão        & \textit{Lack of Cohesion}, v.5 \\
DIT, NOC                  & Herança       & Profundidade; n.\ filhos \\
LOC, LLOC, NOS, NM, NA    & Tamanho       & Linhas, sentenças, métodos, atrib. \\
CC                        & Clones        & \textit{Clone Coverage} \\
CD, CLOC                  & Documentação  & Densidade e linhas de comentário \\
WarningInfo, WarningMajor & Qualidade     & \textit{Warnings} PMD agregados \\
N.\ de bugs               & Alvo          & \textit{Number of bugs} \\
\bottomrule
\end{tabular}
\end{table}

% =====================================================================
\section{Resultados e Discussão}\label{sec:resultados}

\subsection{Medidas de Tendência Central}

A Tabela~\ref{tab:tendencia_class} apresenta média, mediana e moda para as principais variáveis do dataset Class-level.

\begin{table}[htbp]
\caption{Medidas de tendência central -- Class-level}
\label{tab:tendencia_class}
\centering
\footnotesize
\begin{tabular}{lrrr}
\toprule
\textbf{Variável} & \textbf{Média} & \textbf{Mediana} & \textbf{Moda} \\
\midrule
WMC   &  7,39  &  3,00  &  1 \\
CBO   &  5,64  &  3,00  &  2 \\
RFC   &  9,84  &  5,00  &  1 \\
LCOM5 &  1,58  &  1,00  &  1 \\
LOC   & 59,70  & 25,00  &  7 \\
LLOC  & 49,18  & 21,00  &  8 \\
NOS   & 18,77  &  5,00  &  1 \\
NM    &  8,61  &  4,00  &  1 \\
DIT   &  0,98  &  1,00  &  1 \\
CC    &  0,12  &  0,00  &  0 \\
\bottomrule
\end{tabular}
\end{table}

Em todas as métricas de tamanho e complexidade, a média é substancialmente superior à mediana, indicando distribuições fortemente assimétricas à direita. LOC apresenta média de 59,70 e mediana de apenas 25,00. A moda igual a 1 para WMC, RFC, NOS e NM indica que a maioria das classes possui complexidade mínima. DIT apresenta baixa variabilidade (mediana = moda = 1), refletindo hierarquias de herança rasas no Neo4j. Para o File-level, McCC apresenta média 6,06 e mediana 2,00, enquanto LLOC possui média 102,74 e mediana 55,00, confirmando o mesmo padrão.

\subsection{Medidas de Dispersão}

\begin{table}[htbp]
\caption{Medidas de dispersão -- Class-level}
\label{tab:dispersao_class}
\centering
\footnotesize
\begin{tabular}{lrrr}
\toprule
\textbf{Variável} & \textbf{Amplitude} & \textbf{Variância} & \textbf{Desvio Padrão} \\
\midrule
WMC   &     325 &      193,25 &  13,90 \\
CBO   &     161 &       55,57 &   7,45 \\
RFC   &     293 &      256,19 &  16,01 \\
LCOM5 &      71 &       12,33 &   3,51 \\
LOC   &   3.207 &   13.360,81 & 115,59 \\
LLOC  &   2.588 &    8.473,03 &  92,05 \\
NOS   &   1.411 &    2.218,10 &  47,10 \\
NM    &     161 &      183,94 &  13,56 \\
DIT   &       6 &        0,99 &   0,99 \\
CC    &       1 &        0,08 &   0,29 \\
\bottomrule
\end{tabular}
\end{table}

LOC apresenta amplitude de 3.207 e desvio padrão de 115,59 -- quase o dobro da média (59,70) --, confirmando a alta heterogeneidade do tamanho das classes. O coeficiente de variação de LOC é 1,94, o maior entre as métricas de tamanho. DIT, por outro lado, apresenta variância de apenas 0,99, tornando essa métrica praticamente irrelevante como discriminador de \textit{fault-proneness} no Neo4j. No File-level, \textit{Number of developer commits} apresenta variância de 4.106.050,59 e desvio padrão de 2.026,34, decorrente da grande discrepância entre desenvolvedores.

\subsection{Medidas de Posição Relativa e Boxplots}

Para WMC no Class-level: $Q_1 = 1$, $Q_2 = 3$, $Q_3 = 8$, IQR = 7 e $P_{95} = 27$, indicando que 75\% das classes possuem WMC $\leq 8$ enquanto o máximo é 325. Para LOC: $Q_1 = 7$, $Q_2 = 25$, $Q_3 = 60$, IQR = 53. A Fig.~\ref{fig:boxplots_class} evidencia a presença massiva de outliers em todas as métricas. No File-level, McCC ($Q_1 = 1$, $Q_3 = 6$, $P_{90} = 18$) e LLOC ($Q_1 = 20$, $Q_3 = 112$) também exibem assimetria com outliers extremos.

\begin{figure}[htbp]
\centerline{\includegraphics[width=0.95\columnwidth]{figuras/boxplots_class.png}}
\caption{Boxplots das principais métricas -- Class-level.}
\label{fig:boxplots_class}
\end{figure}

\subsection{Avaliação de Outliers}

\begin{table}[htbp]
\caption{Detecção de outliers pelo método IQR -- Class-level}
\label{tab:outliers}
\centering
\footnotesize
\begin{tabular}{lrrr}
\toprule
\textbf{Variável} & \textbf{Outliers} & \textbf{\% Total} & \textbf{Limite Sup.} \\
\midrule
WMC   &   706 &  8,92\% &  18,50 \\
CBO   &   644 &  8,13\% &  14,50 \\
RFC   &   759 &  9,59\% &  24,50 \\
LOC   &   799 & 10,09\% & 140,00 \\
LLOC  &   787 &  9,94\% & 112,50 \\
NOS   &   840 & 10,61\% &  42,00 \\
NM    &   866 & 10,94\% &  19,50 \\
CC    & 1.780 & 22,48\% &   0,00 \\
\bottomrule
\end{tabular}
\end{table}

A proporção de outliers varia entre 5,86\% (NA) e 22,48\% (CC). Para as métricas de tamanho (LOC, LLOC, NOS), aproximadamente 10\% são outliers, consistente com a regra de Pareto em software: uma minoria de componentes concentra a maior complexidade~\cite{nagappan2006}. CC apresenta o maior percentual de outliers porque a maioria dos valores é zero. Os outliers não foram removidos por representarem classes legitimamente complexas do framework.

\subsection{Visualizações Gráficas}

\begin{figure}[htbp]
\centerline{\includegraphics[width=0.95\columnwidth]{figuras/histogramas_class.png}}
\caption{Histogramas das métricas Class-level.}
\label{fig:hist_class}
\end{figure}

Os histogramas (Fig.~\ref{fig:hist_class}) revelam \textbf{assimetria positiva} em \emph{todas} as oito métricas representadas: grande concentração de massa em valores baixos e caudas longas estendendo-se à direita. LOC, LLOC e NOS exibem as caudas mais longas (limites superiores observados de 3.207, 2.588 e 1.411), enquanto WMC, RFC e LCOM5 mostram pico extremo próximo de zero. Essa forma carrega três implicações estatísticas relevantes. \textbf{(a) Tendência central distorcida.} A relação $\bar{x} > \text{mediana} > \text{moda}$ é o marcador clássico de assimetria positiva. \textbf{(b) Inadequação de métodos paramétricos.} Técnicas que pressupõem normalidade tornam-se enviesadas, o que justifica o uso de Spearman e a regressão logística com pesos \textit{cost-sensitive}. \textbf{(c) Compatibilidade com distribuições típicas de software.} A combinação pico-em-zero, decaimento aproximadamente exponencial e cauda pesada é compatível com distribuições \textit{log-normal} ou de \textit{lei de potência}, padrão reportado para métricas em sistemas de grande porte~\cite{radjenovic2013}.

A distribuição de bugs é extremamente desbalanceada: apenas 8 classes com bug (0,10\%) e 6 arquivos com bug (0,14\%) entre milhares de instâncias.

\begin{figure}[htbp]
\centerline{\includegraphics[width=0.75\columnwidth]{figuras/dispersao_loc_wmc.png}}
\caption{Gráfico de dispersão: LOC vs WMC (Class-level).}
\label{fig:disp_loc_wmc}
\end{figure}

A Fig.~\ref{fig:disp_loc_wmc} apresenta a dispersão entre LOC e WMC, evidenciando a correlação forte e positiva entre tamanho e complexidade. As poucas instâncias com bugs (pontos vermelhos) tendem a se posicionar em regiões de valores moderados a altos, embora o número ínfimo de positivos (8 classes com bug) não permita conclusões robustas.

\subsection{Testes de Normalidade}\label{sec:normalidade}

Para todas as variáveis, os testes de Shapiro-Wilk~\cite{shapiro1965} e Lilliefors~\cite{lilliefors1967} rejeitaram a hipótese nula de normalidade ($p < 10^{-50}$). A Tabela~\ref{tab:normalidade} apresenta estatísticas para variáveis representativas.

\begin{table}[htbp]
\caption{Testes de normalidade -- Class-level (amostra de 5.000 obs.\ para SW; seed=42)}
\label{tab:normalidade}
\centering
\footnotesize
\begin{tabular}{lcccc}
\toprule
\textbf{Variável} & \textbf{SW} & \textbf{SW $p$} & \textbf{Lill.\ $D$} & \textbf{Lill.\ $p$} \\
\midrule
WMC & 0,4214 & $< 10^{-83}$ & 0,3081 & $< 2{,}2 \cdot 10^{-16}$ \\
CBO & 0,5874 & $< 10^{-76}$ & 0,2246 & $< 2{,}2 \cdot 10^{-16}$ \\
RFC & 0,5158 & $< 10^{-79}$ & 0,2786 & $< 2{,}2 \cdot 10^{-16}$ \\
LOC & 0,4244 & $< 10^{-83}$ & 0,3113 & $< 2{,}2 \cdot 10^{-16}$ \\
DIT & 0,8267 & $< 10^{-59}$ & 0,2520 & $< 2{,}2 \cdot 10^{-16}$ \\
CC  & 0,4754 & $< 10^{-81}$ & 0,4410 & $< 2{,}2 \cdot 10^{-16}$ \\
\bottomrule
\end{tabular}
\end{table}

A estatística de Shapiro-Wilk varia entre 0,013 (\textit{Number of bugs}) e 0,827 (DIT), sempre muito distante de 1. A escolha pelo teste de Lilliefors~\cite{lilliefors1967} em detrimento do tradicional Kolmogorov-Smirnov (K-S) decorre de uma restrição técnica do K-S clássico: sua tabela de valores críticos só é corretamente calibrada quando a média ($\mu$) e o desvio padrão ($\sigma$) da distribuição teórica de referência são fixados \textit{a priori}. Como $\mu$ e $\sigma$ precisam ser estimados a partir da própria amostra, o K-S com a tabela padrão torna-se excessivamente conservador. O teste de Lilliefors corrige essa distorção recalculando os valores críticos de modo a respeitar o nível de significância nominal quando $\mu$ e $\sigma$ são estimados amostralmente. A convergência dos dois testes com $p < 10^{-50}$ elimina qualquer dúvida sobre a violação de normalidade.

\subsection{Análise de Correlação}

Dado que todas as variáveis violam normalidade, a correlação de Pearson foi complementada com a de Spearman, mais robusta a \textit{outliers} e distribuições assimétricas.

\begin{figure}[htbp]
\centerline{\includegraphics[width=0.85\columnwidth]{figuras/correlacao_heatmap_class.png}}
\caption{Correlação de Pearson -- Class-level.}
\label{fig:corr_class}
\end{figure}

A Fig.~\ref{fig:corr_class} apresenta a matriz de correlação de Pearson. \textbf{Correlações fortes} ($|r| \geq 0{,}7$, Pearson): LOC x LLOC: $r = 0{,}99$; LOC x NOS: $r = 0{,}95$; WMC x LOC: $r = 0{,}89$; WMC x RFC: $r = 0{,}82$; CBO x RFC: $r = 0{,}80$. \textbf{Moderadas} ($0{,}4 \leq |r| < 0{,}7$): WMC x CBO: $r = 0{,}51$; WMC x NM: $r = 0{,}53$; CBO x LOC: $r = 0{,}59$. As correlações de Spearman (Fig.~\ref{fig:corr_spearman_class}) apresentam padrão semelhante com magnitudes geralmente mais elevadas (e.g., LOC x LLOC: $\rho = 0{,}99$; NOS x NL: $\rho = 0{,}72$). Para o dataset File-level, as correlações mais fortes ocorrem entre McCC x LLOC ($r = 0{,}72$) e \textit{N.\ previous modifications} x \textit{N.\ committers} ($r = 0{,}77$; Spearman: $\rho = 0{,}90$).

\begin{figure}[htbp]
\centerline{\includegraphics[width=0.85\columnwidth]{figuras/correlacao_spearman_class.png}}
\caption{Correlação de Spearman -- Class-level.}
\label{fig:corr_spearman_class}
\end{figure}

\subsection{Análise de Multicolinearidade (VIF)}

O \textit{Variance Inflation Factor} (VIF) foi calculado para quantificar formalmente a multicolinearidade entre os preditores~\cite{montgomery2014}. VIF $> 10$ indica multicolinearidade severa.

\begin{figure}[htbp]
\centerline{\includegraphics[width=0.9\columnwidth]{figuras/vif_barplot.png}}
\caption{VIF dos preditores da regressão logística (antes da remoção).}
\label{fig:vif}
\end{figure}

No Class-level, RFC (VIF = 20,18), WMC (15,17) e LOC (10,98) excederam o limiar de 10 (Fig.~\ref{fig:vif}). Após remoção iterativa de RFC, os VIF finais ficaram abaixo de 10 (WMC: 8,75; LOC: 9,16; demais $\leq$ 5). No File-level, todos os VIF foram aceitáveis ($\leq$ 5,58). A remoção de RFC do modelo Class-level é justificada pela sua redundância com WMC ($r = 0{,}82$) e CBO ($r = 0{,}80$).

\subsection{Modelagem Estatística -- Regressão Logística}

Foi adotada a regressão logística binária com pesos \textit{cost-sensitive} inversamente proporcionais à frequência de cada classe ($w_+ = N / (2 \cdot N_+)$; $w_- = N / (2 \cdot N_-)$), evitando classificação trivial como classe majoritária.

\subsubsection{Modelo Class-level}
O modelo utilizou os preditores remanescentes após remoção de VIF $> 10$: WMC, CBO, LOC, LCOM5, NL, DIT, CC. Os coeficientes apresentaram magnitude extrema (ordem $10^{12}$--$10^{15}$), indicando que o algoritmo \textbf{não convergiu} -- característica de separação completa dos dados, sintoma do desbalanceamento extremo (8 classes com bug em 7.917, EPV $\approx 1$). O Pseudo $R^2$ (McFadden) foi $-10{,}76$. Valores entre 0,2 e 0,4 indicam ajuste excelente~\cite{mcfadden1974}.

\subsubsection{Modelo File-level}
Os preditores foram: McCC, LLOC, N.\ previous modifications, N.\ previous fixes, N.\ committers e N.\ developer commits. O Pseudo $R^2$ (McFadden) foi 0,59, porém o algoritmo também não convergiu. Apenas LLOC não foi significativo ($p = 0{,}30$); os demais resultados devem ser interpretados com cautela dado o número ínfimo de positivos (6 arquivos com bug).

\subsubsection{Métricas de Avaliação e Validação Cruzada}
Para evitar \textit{overfitting}, foi implementada validação cruzada estratificada com 10 \textit{folds}. O \textit{threshold} ótimo foi determinado pelo Youden Index ($J = \text{Sensit.} + \text{Espec.} - 1$).

\begin{table}[htbp]
\caption{Métricas de avaliação: treino completo e 10-fold CV}
\label{tab:metricas_avaliacao}
\centering
\footnotesize
\begin{tabular}{lrrrr}
\toprule
 & \multicolumn{2}{c}{\textbf{Treino completo}} & \multicolumn{2}{c}{\textbf{10-fold CV}} \\
\cmidrule(lr){2-3} \cmidrule(lr){4-5}
\textbf{Métrica} & \textbf{Class} & \textbf{File} & \textbf{Class} & \textbf{File} \\
\midrule
AUC-ROC     & 0,674 & 0,968 & 0,660 & 0,805 \\
AUC-PR      & 0,007 & 0,020 & ---   & ---   \\
Sensitivity & 0,375 & 1,000 & 0,625 & 0,833 \\
Precision   & 0,014 & 0,024 & 0,002 & 0,007 \\
F1-score    & 0,027 & 0,047 & 0,004 & 0,014 \\
Bal.\ Acc.  & 0,674 & 0,972 & 0,664 & 0,836 \\
MCC         & 0,069 & 0,151 & 0,023 & 0,068 \\
\bottomrule
\end{tabular}
\end{table}

\begin{figure}[htbp]
\centerline{\includegraphics[width=0.95\columnwidth]{figuras/curva_roc.png}}
\caption{Curvas ROC com Youden Index para Class-level e File-level.}
\label{fig:roc}
\end{figure}

A Fig.~\ref{fig:roc} apresenta as curvas ROC com os pontos ótimos pelo Youden Index. A degradação na validação cruzada confirma \textit{overfitting}, especialmente no File-level (AUC: 0,97 $\rightarrow$ 0,81). Os valores de MCC próximos de zero (0,023 e 0,068 no CV) indicam poder preditivo praticamente nulo. A ``acurácia'' de 97--99\% é enganosa: um classificador trivial que sempre prediz ``sem bug'' obtém desempenho similar. As curvas Precision-Recall (AUC-PR $\approx 0{,}007$ e $0{,}020$) confirmam a inadequação dos dados para predição.

\subsection{Discussão}

As métricas do Neo4j seguem distribuições fortemente assimétricas, consistente com a literatura~\cite{radjenovic2013,gyimothy2005}. As correlações fortes entre LOC, LLOC, NOS, WMC e RFC indicam dimensões sobrepostas de tamanho e complexidade, confirmadas pela análise VIF. A correlação extremamente fraca entre métricas e bugs (máximo $r = 0{,}06$) contrasta com estudos anteriores~\cite{basili1996,gyimothy2005}, explicável pelo desbalanceamento extremo e provável subreporte de defeitos. Quanto às métricas CK: DIT e NOC possuem variância quase zero no Neo4j (design ``flat''). As métricas WMC e CBO, apontadas como melhores preditores na literatura~\cite{radjenovic2013}, também não apresentaram poder preditivo aqui, o que reforça que a limitação é do dataset, não das métricas.

\subsection{Ameaças à Validade}\label{sec:ameacas}

\textbf{Validade interna.} (i) Desbalanceamento extremo (8 classes e 6 arquivos com bug; EPV $\approx 1$, abaixo do mínimo de 10~\cite{hosmer2013}), responsável pela separação completa e não-convergência; (ii) K-S com parâmetros estimados seria conservador -- corrigido com Lilliefors; (iii) Multicolinearidade severa (LOC x LLOC: $r = 0{,}99$) -- corrigida via remoção iterativa de variáveis com VIF $> 10$; (iv) Avaliação no próprio conjunto de treino -- corrigida com 10-fold CV estratificada; (v) Subamostragem aleatória para Shapiro-Wilk com semente fixa (\texttt{seed=42}) -- mitigada pela convergência com Lilliefors.

\textbf{Validade externa.} (i) Resultados restritos a um único projeto (Neo4j) em uma única versão; (ii) domínio específico (banco de dados em grafo) com arquitetura plana -- a baixa variância de DIT e NOC é provavelmente idiossincrasia; (iii) viés de captura do GitHub Bug Dataset~\cite{toth2016}: apenas defeitos vinculados a \textit{commits} de correção são registrados.

\textbf{Validade de construto.} (i) Métricas CK são \textit{proxies} estruturais imperfeitos da qualidade real; (ii) \textit{Number of bugs} operacionaliza ``defeito'' como ``bug historicamente registrado''; (iii) LCOM5 possui múltiplas formulações; (iv) a agregação de regras PMD em \texttt{WarningInfo}/\texttt{WarningMajor} sacrifica resolução semântica.

\textbf{Validade de conclusão.} (i) Baixíssimo poder estatístico para inferências sobre a variável-alvo; (ii) cerca de 170 correlações pareadas no Class-level sem correção formal para comparações múltiplas -- as associações fortes ($|r| \geq 0{,}7$) coincidem em Pearson e Spearman; (iii) os $p$-valores extremamente pequenos refletem a sensibilidade dos testes a $n$ grande; (iv) as convenções internas do SourceMeter afetam todas as métricas sistematicamente.

% =====================================================================
\section{Conclusão}\label{sec:conclusao}

Este trabalho apresentou uma análise estatística rigorosa de métricas de código-fonte do framework Neo4j (GitHub Bug Dataset v1.1), com testes de normalidade corrigidos (Lilliefors), correlações de Pearson e Spearman, análise VIF, regressão \textit{cost-sensitive} e validação cruzada.

Os principais achados foram: (1) todas as métricas seguem distribuições assimétricas com correlações fortes entre tamanho e complexidade (LOC x LLOC: $r = 0{,}99$); (2) nenhuma variável é normal (todos os $p < 10^{-50}$); (3) a proporção de outliers varia entre 5\% e 22\%; (4) o \textbf{principal achado é a inadequação do dataset para predição de defeitos}: com EPV $\approx 1$, os modelos são não confiáveis e o MCC próximo de zero na CV (0,023 e 0,068) confirma poder preditivo nulo.

Como trabalhos futuros: (i) combinar múltiplas versões do Neo4j para aumentar instâncias positivas e viabilizar EPV $\geq 10$; (ii) aplicar \textit{Random Forest} ou \textit{gradient boosting}, que tratam desbalanceamento nativamente; (iii) explorar \textit{cross-project defect prediction}.

% =====================================================================
\begin{thebibliography}{00}
\bibitem{pressman2014} R.~S. Pressman and B.~R. Maxim, \textit{Software Engineering: A Practitioner's Approach}, 8th ed. New York: McGraw-Hill Education, 2014.
\bibitem{chidamber1994} S.~R. Chidamber and C.~F. Kemerer, ``A metrics suite for object oriented design,'' \textit{IEEE Trans. Softw. Eng.}, vol.~20, no.~6, pp.~476--493, 1994.
\bibitem{gyimothy2005} T. Gyim{\'o}thy, R. Ferenc, and I. Siket, ``Empirical validation of object-oriented metrics on open source software for fault prediction,'' \textit{IEEE Trans. Softw. Eng.}, vol.~31, no.~10, pp.~897--910, 2005.
\bibitem{basili1996} V.~R. Basili, L.~C. Briand, and W.~L. Melo, ``A validation of object-oriented design metrics as quality indicators,'' \textit{IEEE Trans. Softw. Eng.}, vol.~22, no.~10, pp.~751--761, 1996.
\bibitem{toth2016} Z. T{\'o}th, P. Gyim{\'e}si, and R. Ferenc, ``A public bug database of GitHub projects and its application in bug prediction,'' in \textit{Proc. Int. Conf. Comput. Sci. Appl. (ICCSA)}, 2016, pp.~625--638.
\bibitem{radjenovic2013} D. Radjenovi{\'c}, M. Heriko, R. Torkar, and A. {\v{Z}}ivkovi{\v{c}}, ``Software fault prediction metrics: A systematic literature review,'' \textit{Inf. Softw. Technol.}, vol.~55, no.~8, pp.~1397--1418, 2013.
\bibitem{montgomery2014} D.~C. Montgomery and G.~C. Runger, \textit{Applied Statistics and Probability for Engineers}, 6th ed. Hoboken, NJ: John Wiley \& Sons, 2014.
\bibitem{shapiro1965} S.~S. Shapiro and M.~B. Wilk, ``An analysis of variance test for normality (complete samples),'' \textit{Biometrika}, vol.~52, no.~3/4, pp.~591--611, 1965.
\bibitem{nagappan2006} N. Nagappan, T. Ball, and A. Zeller, ``Mining metrics to predict component failures,'' in \textit{Proc. 28th Int. Conf. Softw. Eng. (ICSE)}, 2006, pp.~452--461.
\bibitem{lilliefors1967} H.~W. Lilliefors, ``On the Kolmogorov-Smirnov test for normality with mean and variance unknown,'' \textit{J. Amer. Statist. Assoc.}, vol.~62, no.~318, pp.~399--402, 1967.
\bibitem{mcfadden1974} D. McFadden, ``Conditional logit analysis of qualitative choice behavior,'' in \textit{Frontiers in Econometrics}, P. Zarembka, Ed. New York: Academic Press, 1974, pp.~105--142.
\bibitem{hosmer2013} D.~W. Hosmer, S. Lemeshow, and R.~X. Sturdivant, \textit{Applied Logistic Regression}, 3rd ed. Hoboken, NJ: John Wiley \& Sons, 2013.
\end{thebibliography}

\end{document}
